\begin{abstract}
Conditional hierarchical classifiers can reflect broad-to-specific diagnostic
organisation, but an incorrect early decision may prevent the appropriate
downstream classifier from receiving an image. This work evaluates that
routing dependency using a two-stage EfficientNet-B0 hierarchy and two flat
comparators on the same leakage-aware, locked 3,668-image internal-test cohort
derived from ISIC 2019. The endpoint classes are non-malignant, melanoma,
basal cell carcinoma, and squamous cell carcinoma.

Flat DenseNet-121 produced the highest descriptive point estimates for
accuracy (0.7914), macro-F1 (0.6351), and weighted-F1 (0.7862), whereas flat
EfficientNet-B0 achieved the highest balanced accuracy (0.6503). The
predicted-gate hierarchy obtained 0.7402 accuracy and 0.6054 macro-F1. Paired
analyses supported DenseNet-121's accuracy and weighted-F1 gains
over both EfficientNet-B0 systems, while its macro-F1 advantage was
supported only against the hierarchy. Its balanced accuracy was lower than
flat EfficientNet-B0, while the difference from the hierarchy was
inconclusive. In the backbone-matched
EfficientNet-B0 comparison, the flat-minus-hierarchy macro-F1 difference was
0.0139 with a paired-bootstrap 95\% confidence interval of
$[-0.0143,\,0.0420]$. Exact McNemar testing yielded $p=0.8207$, detecting no
statistically significant difference.

Replacing the learned route with the ground-truth route increased hierarchical
macro-F1 to 0.7937, corresponding to a routing-associated loss of 0.1883.
Stage~1 blocked 20.079\% of malignant images before subtype classification.
DenseNet-121 nevertheless achieved only 0.2979 SCC recall. Similar aggregate
results can therefore conceal substantial routing and minority-class failure.
\end{abstract}

\begin{IEEEkeywords}
dermoscopy, skin lesion classification, hierarchical classification,
conditional inference, routing errors, EfficientNet, DenseNet, class imbalance
\end{IEEEkeywords}
